\section{Theory}
\label{sec:theory}
\subsection{Cost functions and their gradients}
The loss function is the most important function 
for training, since training a model
means to minimize the cost function of the model.  
To minimize the cost function, we use backpropagation
to find the gradient of the cost, with regards
to  and gradient descent.

\subsubsection{MSE}
When we have our calculated output value $\tilde{y}$ and the true
output value $y$, with n data points, the mean-squared error and its derivative is
\begin{equation}
    MSE(y,\tilde{y}) = \frac{1}{n}\sum_{i=0}^{n-1}(y - \tilde{y})^2 \quad 
    \frac{1}{n}||y - X\Theta||_2^2
    \label{eq: MSE}
\end{equation}
\paragraph{L2 norm} which corresponds to the Ridge cost function.
\begin{equation}
    MSE + \lambda||\Theta||_2^2 \quad \Theta(W, b)
    \label{eq: RidgeCost}
\end{equation}
\paragraph{L1 norm} which corresponds to the Lasso cost function.
\begin{equation}
    MSE + \lambda||\Theta|| \quad \Theta(W, b)
    \label{eq: LassoCost}
\end{equation}
\cite[p.~86-91]{LectureNotes}

\subsubsection{Binary cross entropy: log loss}
With two classes, we assume a binary output $y \in \{0, 1\}$, and a predicted probability $p \in [0, 1]$ which is
the output of the \hyperref[eq:Sigmoid]{sigmoid function} of $z_i$. 
Where $z_i = x_i^T\Theta$ is the linear combination of the input features and the model parameters $\Theta$. 
For a single example, data point $i$, the binary cross entropy cost function is defined as 
\begin{align}
    C(\Theta) &= -\sum_{i=1}^{n}(y_i \log p_i + 
    (1 - y_i)\log[1 - p_i]) \\
    p_i &= \sigma(z_i) = \frac{1}{1 + \exp(-z_i)} \\
    \label{eq: BinaryCost}
\end{align}
The cost function is derived from the maximum likelihood estimation of the Bernoulli distribution,
the negative log likelihood of the Bernoulli distribution. This is known as the cross entropy, and
it is the loss function used to train binary classification models. 
With the linear function $z_i = x_i^T\Theta$, the gradient of the cost function with respect 
to the model parameters $\Theta$ is given by
\begin{equation}
    \frac{dC(\theta)}{d\theta} = -X^T(y - p)
    \label{eq: BinaryCostGradient}
\end{equation}
For N data points, the mean binary cross entropy is defines as $$\frac{1}{N}\sum_{i=1}^{N}C(\theta)_i$$
And the term $1/N$ follows through te calculations to give the mean binary cross entropy gradient, which is given by
\begin{equation}
    \nabla_\Theta C(\Theta) = -\frac{1}{n}X^T(y - p)
    \label{eq:MeanBinaryCostGradient}
\end{equation}
The cross entropy L1 and L2 norm is supplemented like the ordinary MSE. 

\subsubsection{Multiclass cross entropy cost function}
For k klasses, one uses the soft max function to calculate the probability of each class.
Again we have the linear function $z_k = X^T\Theta_k$ for each class, over all features and data points, 
the softmax functions is defined as
\begin{equation}
    p(C = k \mid x) = \frac{\exp(z_k)}
                 {\sum_{l=0}^{K-1}\exp(z_l)} \\
        k = 0, \ldots, K-1
    \label{eq:SoftmaxFunc}
\end{equation}
The nominator is for class k, and the denominator is the sum of all classes.
The multiclass cross entropy cost function is defined as
\begin{equation}
    C(\Theta) = -\sum_{i=1}^{n}\sum_{k=0}^{K-1}y_{ik}\log p(C=k|x_i)
    \label{eq: SoftmaxCost}
\end{equation}
Which is summed over all data points and classes.
The gradient of the multiclass cross entropy cost function is given by
\begin{equation}
    \frac{dC}{d\Theta_k} = -X^T(y_k - p_k)
    \label{eq: SoftmaxCostGradient}
\end{equation}
The cross entropy L1 and L2 norm is supplemented like the ordinary MSE.  
Stability:
The softmax function can be numerically unstable for large values of $z_k$, for a large value the e
exponential function can overflow. Since the softmax is invariant udner subtracting a constant frome
every score, one subtracts the largest before exponentiating:
$$p_k = \frac{\exp(z_k - z_{max})}{\sum_{l}\exp(z_l - z_{max})}$$, $z_k = x^T\Theta_k$. 
\cite[p.~151-156]{LectureNotes}

\subsection{Activation functions and their gradients}
% Equations for the activation functions and their gradients.
% LLM Declaration: Equations generated with Codex assistance.
\subsubsection{Sigmoid logit function}

\begin{equation}
    \sigma(z) = \frac{1}{1 + \exp(-z)}
    \label{eq:Sigmoid}
\end{equation}
\begin{equation}
    \sigma'(z) = \sigma(z)(1 - \sigma(z))
    \label{eq : SigmaPrime}
\end{equation}

\subsubsection{RELU function}

\begin{equation}
    ReLU(z) = \max(0, z)
    \label{eq : ReLU}
\end{equation}

\begin{equation}
    ReLU'(z) = \begin{cases}
        0 & z < 0 \\
        1 & z > 0
    \end{cases}
    \label{eq : ReLUPrime}
\end{equation}
\cite[p.~214-242]{LectureNotes}


